% ============================================================
%  Appendix B: Dataset Statistics
% ============================================================

\chapter{Dataset Statistics and Characteristics}
\label{app:datasets}

This appendix provides detailed statistical characteristics for the three datasets used
across the empirical studies in this thesis: RanSAP, RDset, and RansKnow.

% ------------------------------------------------------------
\section*{B.1\quad RanSAP — Runtime Behavioural Dataset}
% ------------------------------------------------------------

The RanSAP dataset \cite{hirano2022ransap} captures system API call behaviour recorded
during controlled sandbox execution of ransomware and benign software samples.
Samples are represented as 30-second time-window API trace feature vectors capturing
I/O access patterns, file enumeration sequences, and encryption-indicative operations.
Table~\ref{tab:app_ransap} summarises the class and split distribution used in this thesis.

\begin{table}[ht]
\centering
\caption{RanSAP dataset class distribution and train/test split.}
\label{tab:app_ransap}
\begin{tabular}{@{}lrrr@{}}
\toprule
\textbf{Class} & \textbf{Training} & \textbf{Test} & \textbf{Total} \\
\midrule
Ransomware      & ---  & ---  & --- \\
Benign          & ---  & ---  & --- \\
\midrule
\textbf{Total}  & ---  & ---  & --- \\
\bottomrule
\end{tabular}
\end{table}

\noindent\textit{Note:} Exact split figures are reported within each empirical chapter.
The split ratio used throughout is 80\% training, 10\% validation, 10\% test, with
stratified sampling to preserve class proportions.

\begin{table}[ht]
\centering
\caption{RanSAP feature group summary.}
\label{tab:app_ransap_features}
\begin{tabular}{@{}lll@{}}
\toprule
\textbf{Feature Group} & \textbf{Count} & \textbf{Description} \\
\midrule
File I/O API calls        & ---  & Read, write, create, delete, rename operations \\
Registry operations       & ---  & Key creation, modification, deletion events \\
Network API calls         & ---  & Socket creation, DNS resolution, HTTP requests \\
Process API calls         & ---  & Process and thread creation, injection events \\
Cryptographic API calls   & ---  & Encryption/decryption routine invocations \\
\midrule
\textbf{Total features}   & 16   & Per 30-second time window \\
\bottomrule
\end{tabular}
\end{table}

% ------------------------------------------------------------
\section*{B.2\quad RDset — Static Semantic Feature Dataset}
% ------------------------------------------------------------

The RDset provides PE-format static analysis features extracted from portable executable
file headers without dynamic execution, enabling rapid pre-execution screening
\cite{GogineniDV22ForeseerRanSAP}.
Features encode the structural and semantic properties of executable files through entropy
profiles, import address table statistics, and PE header field values.

\begin{table}[ht]
\centering
\caption{RDset class distribution and train/test split.}
\label{tab:app_rdset}
\begin{tabular}{@{}lrrr@{}}
\toprule
\textbf{Class} & \textbf{Training} & \textbf{Test} & \textbf{Total} \\
\midrule
Ransomware      & ---  & ---  & --- \\
Benign          & ---  & ---  & --- \\
\midrule
\textbf{Total}  & ---  & ---  & --- \\
\bottomrule
\end{tabular}
\end{table}

\begin{table}[ht]
\centering
\caption{RDset feature group summary.}
\label{tab:app_rdset_features}
\begin{tabular}{@{}lll@{}}
\toprule
\textbf{Feature Group} & \textbf{Count} & \textbf{Description} \\
\midrule
Section entropy values   & ---  & Per-section byte entropy (Shannon entropy) \\
Import address table     & ---  & DLL names and imported function counts \\
PE header fields         & ---  & Machine type, subsystem, characteristics flags \\
Positional embeddings    & 384  & Final input representation after embedding \\
\midrule
\textbf{Total (raw)}     & ---  & Before positional embedding \\
\bottomrule
\end{tabular}
\end{table}

% ------------------------------------------------------------
\section*{B.3\quad RansKnow — Threat-Intelligence Dataset}
% ------------------------------------------------------------

RansKnow is a researcher-constructed dataset sourced from 307 publicly available YouTube
security briefings, conference presentations, and practitioner DFIR walkthroughs.
Each sample encodes the observed MITRE ATT\&CK tactic profile and associated tooling of
a documented ransomware incident or benign activity report.
Binary labels are assigned by the threshold rule:
\textit{Tactic\_Total} $>$ 10 \textbf{OR} \textit{Tool\_Total} $>$ 2.

\begin{table}[ht]
\centering
\caption{RansKnow class distribution.}
\label{tab:app_ransknow_dist}
\begin{tabular}{@{}lrr@{}}
\toprule
\textbf{Class} & \textbf{Samples} & \textbf{Percentage} \\
\midrule
Ransomware (positive)  & ---  & --- \\
Benign (negative)      & ---  & --- \\
\midrule
\textbf{Total}         & 307  & 100\% \\
\bottomrule
\end{tabular}
\end{table}

\begin{table}[ht]
\centering
\caption{RansKnow full feature specification (23 features).}
\label{tab:app_ransknow_features}
\renewcommand{\arraystretch}{1.3}
\begin{tabular}{@{}llp{7.5cm}@{}}
\toprule
\textbf{Feature Group} & \textbf{No.} & \textbf{Description} \\
\midrule
MITRE ATT\&CK Tactics & 10 & Binary/count indicators for Initial Access, Execution, Persistence, Privilege Escalation, Defence Evasion, Credential Access, Discovery, Lateral Movement, Collection, and Exfiltration \\
\addlinespace
Tool Counts & 8 & Occurrence counts for tools commonly observed in ransomware chains: Cobalt Strike, Mimikatz, PsExec, PowerShell Empire, Rclone, MEGAsync, WinSCP, AnyDesk \\
\addlinespace
Platform Signals & 3 & Indicators for Windows, Linux, and ESXi/VMware as primary target platform \\
\addlinespace
Aggregate Totals & 2 & \textit{Tactic\_Total} (sum of tactic indicators) and \textit{Tool\_Total} (sum of tool counts) \\
\midrule
\textbf{Total} & 23 & \\
\bottomrule
\end{tabular}
\end{table}

\noindent The labelling threshold (\textit{Tactic\_Total} $>$ 10 OR \textit{Tool\_Total}
$>$ 2) was established through examination of practitioner-documented incident reports
and reflects operational consensus on the minimum tactical breadth and tooling complexity
distinguishing a ransomware deployment from benign or low-sophistication activity.
